\documentclass[12pt]{article}
\usepackage[utf8]{inputenc}
\usepackage[T1]{fontenc}
\usepackage[margin=1in]{geometry}
\usepackage{setspace}
\usepackage{xcolor}
\usepackage[colorlinks=true, urlcolor=blue!60!black, linkcolor=blue!60!black]{hyperref}
\setlength{\parindent}{0em}
\setlength{\parskip}{0.45em}
\renewcommand{\baselinestretch}{1.0}
\begin{document}
\begin{center}
{\Large \textit{Massachusetts v.\ EPA}, 549 U.S.\ 497 (2007)\par}
\vspace{0.3em}
{\normalsize Case Brief --- One-Page Edition\par}
\end{center}

\textbf{Memory Jogger.} With ``special solicitude'' for States, Massachusetts could challenge EPA's refusal to regulate vehicle greenhouse gases, which are ``air pollutants'' (\href{https://supreme.justia.com/cases/federal/us/549/497/}{Justia}).

\textbf{Facts.} EPA denied a petition to regulate greenhouse gases from new cars, saying the Clean Air Act did not reach climate change and regulation was unwise. Massachusetts showed rising seas eroding its state-owned coast.

\textbf{Procedural History.} A fractured D.C.\ Circuit panel denied review.

\textbf{Votes.} Reversed 5--4: Stevens for the Court; Roberts (standing) and Scalia (merits) dissented, with Thomas and Alito.

\textbf{Holding.} A State asserting a statutory procedural right to protect its quasi-sovereign interests gets special solicitude; Massachusetts showed injury, causation, and incremental redressability. Greenhouse gases are ``air pollutants,'' and EPA may not refuse to act for extrastatutory reasons.

\textbf{Reasoning.} States are ``not normal litigants,'' and agencies may ``whittle away'' at big problems; the Act's repeated ``any'' is unambiguous. Roberts called global warming ``harmful to humanity at large'' and the ruling ``SCRAP for a new generation''; Scalia would have deferred to EPA.

\textbf{Examples.} \emph{Same-Side:} a State protecting its tidal lands and fisheries. \emph{Opposite-Side:} a private group alone alleging future climate harm. \emph{Fence-Sitter:} a State claiming lost revenue from deregulation.

\textbf{Critique.} ``Special solicitude'' lacks clear constitutional basis, was never defined, and spawned partisan state litigation.

\textbf{Key Quotation.} ``Given that procedural right and Massachusetts' stake in protecting its quasi-sovereign interests, the Commonwealth is entitled to special solicitude in our standing analysis.'' --- Justice Stevens, 549 U.S.\ at 520.

\end{document}
